% 卒業論文ドラフト
% ビルド: latexmk -lualatex main.tex
% 研究室のテンプレートが判明したら documentclass とタイトル部を差し替える
\documentclass[report, jafontscale=1.0]{jlreq}

\usepackage{luatexja-fontspec}
\usepackage{graphicx}
\usepackage{booktabs}
\usepackage{amsmath}
\usepackage{amssymb}
\usepackage{url}
\usepackage[hidelinks]{hyperref}

\newcommand{\todo}[1]{\par\noindent\textbf{[TODO: #1]}\par}

\title{極大クリーク列挙の並列化\\
  \large Eppstein アルゴリズムの最外層並列化と性能限界の実測分析（仮題）}
\author{中島 蒼}
\date{2027年2月（提出予定）}

\begin{document}

\maketitle

\begin{abstract}
% 概要は最後に書く。本文の各章が固まってから、1段落×4（背景・手法・結果・結論）で構成する。
\todo{概要（最後に書く）}
\end{abstract}

\tableofcontents

\chapter{はじめに}
\label{ch:introduction}

\section{背景}
\label{sec:intro-background}

グラフのなかで互いに全て隣接している頂点の集合を\textbf{クリーク}という。
どの頂点を足してもクリークでなくなるものを\textbf{極大クリーク}と呼び、グラフに含まれる極大クリークをすべて列挙する問題を\textbf{極大クリーク列挙}という。
極大クリークは、要素どうしが密につながった集団をそのまま取り出す構造であり、ネットワークとして表せるデータから結びつきの強いまとまりを見つける道具として使われる。

こうしたまとまりの抽出は、データ分析の一部として現れる。
扱うデータは年々大きくなり、その分析を速く回せるかどうかが、データを持つ側がそこから何を引き出せるかを左右する。
分析の一段として極大クリーク列挙を走らせるなら、その列挙自体が遅ければ後続の作業も待たされる。
本研究は、この列挙を高速化する手立てとして並列化を扱う。
たとえば、極大クリークを種として複雑ネットワークから重複しあうコミュニティを抽出する手法があり、共著ネットワークやタンパク質相互作用ネットワークに適用されている \cite{palla2005}。

\section{問題}
\label{sec:intro-problem}

極大クリーク列挙は、出力の量に敏感な問題である。
$n$ 頂点のグラフが持ちうる極大クリークの数は最悪で $3^{n/3}$ に達し、この上限は Moon と Moser の極値グラフが実際に達成する \cite{moon1965}。
出力そのものが指数的に膨らみうるため、入力サイズが同じでも、列挙にかかる時間はグラフの形によって大きく変わる。

大きなグラフでは、逐次実行の時間が実用上の壁になる。
本研究で主対象とした soc-sign-epinions（頂点数 131{,}828、極大クリーク数 約2{,}223万）では、逐次実行に、縮退順序を用いる Eppstein のアルゴリズム \cite{eppstein2010,eppstein2013} で約95秒、ピボット選択を用いる CLIQUES \cite{tomita2006} で約167秒を要した。
一度の分析で何度も列挙を回す場面を考えると、この所要時間は無視できない。
手元の計算機が複数のコアを備える以上、それらを使って列挙を速くできないかを問うのは自然である。

\section{本研究のアプローチ}
\label{sec:intro-approach}

本研究は、Eppstein のアルゴリズムの最外層を並列化する。
Eppstein のアルゴリズムは、グラフの疎さを表す\textbf{縮退数} $d$ を用いて最悪計算量を $O(d\,n\,3^{d/3})$ に抑える変種であり \cite{eppstein2010}、最外層を縮退順序に沿って回る。
この最外層の各頂点を根とする部分問題は互いに独立しているため、頂点ごとの部分問題を複数の worker へ配れば、そのまま並列化できる。

手法の選び方ではなく、進め方に本研究の重点を置く。
最外層を素朴に並列化した実装から出発し、性能を測り、伸びを止めている要因を一つ特定しては潰す、という実測駆動の反復で進めた。
律速要因を机上で見積もって一度に設計するのではなく、測定が示した順に、プロセス起動の固定費、worker 間の負荷の偏り、計算機の非対称なコア構成、並列化に伴うメモリのコピーを、順に切り分けていった。
測定は単発の値の振れが大きかったため、同一条件を3回計測して最小値を採る方針で揃えている。

\section{本研究の貢献}
\label{sec:intro-contributions}

本研究の貢献は次の四点である。

\begin{itemize}
  \item \textbf{逐次実装での二手法の実測比較}：次数分布と疎密の異なる10種類のグラフで、CLIQUES と Eppstein のアルゴリズムを同一条件で比較した。両者の優劣は縮退数のグラフサイズに対する比 $d/n$ でほぼ説明でき、疎で大きいグラフほど Eppstein が有利になる（最大で木の16.6倍）一方、密なグラフでは互角に収束し、明確に負けるのは完全グラフだけだった。この結果が、疎で大きいグラフを並列化の主対象に選ぶ根拠になっている。
  \item \textbf{律速要因の測定による切り分け}：素朴な並列実装が soc-sign-epinions で頭打ちになる原因を、フェーズ別の時間分解と頂点ごとの負荷計測で特定した。プロセスプールの起動固定費が worker 数にほぼ比例して伸びること（spawn 方式の8 worker で約1.27秒）、仕事が縮退順序の末尾の少数の重い頂点に偏ること（上位1\%の頂点が全所要時間の94.7\%を占める）、そして計算機の非対称なコア構成が効いていることを、いずれも測定値で示した。
  \item \textbf{静的な負荷分散の効果と限界}：重い頂点と軽い頂点を各バッチに混ぜる配分（interleave）で、worker の遊休を32.6\%から0.1\%までほぼ消せることを確かめた。それにもかかわらず実行時間の短縮は4.3\%にとどまり、その原因が、性能コアと効率コアからなる非対称なコア構成にあることを、コアを固定した計測で裏づけた（効率コア固定時の速度比は性能コア比0.26。これは計測方法上の下限値で、全力時は約0.43と推定される）。遊休を測る指標がコアの均質性を暗黙に仮定していたことを、この計測が明らかにした。
  \item \textbf{起動固定費の削減}：プロセス起動方式を spawn から fork に変えるだけで、worker 数に比例していた起動固定費をほぼ消せることを示した（soc-sign-epinions の8 worker で起動固定費が約0.03秒）。これにより、素朴な実装では並列化が損だった小さいグラフのうち4種が、4 worker で約2倍の高速化に転じた。共有メモリを併用する変種と合わせ、soc-sign-epinions での8 worker の実行時間を約26.5秒（逐次の94.9秒に対し約3.6倍）まで短縮した。
\end{itemize}

\section{本論文の構成}
\label{sec:intro-outline}

第2章では、クリーク、縮退数、Bron--Kerbosch 系のアルゴリズムなど、本論文で用いる用語と対象アルゴリズムを整理する。
第3章では、極大クリーク列挙の並列化に関する先行研究を、動的スケジューリングと共有メモリ化の観点から概観する \cite{schmidt2009,das2018parmce}。
第4章では、最外層並列化の素朴な実装と、測定を経て加えた改良を述べる。
第5章では、逐次比較と並列化の各段階の測定結果を示し、律速要因の切り分けを報告する。
第6章で、得られた知見と今後の課題をまとめる。

\chapter{準備}
\label{ch:preliminaries}

本章では、以降の章で極大クリーク列挙アルゴリズムを並列化するために必要な概念を定義する。
グラフとクリークの基本用語から始め、逐次アルゴリズムである Bron--Kerbosch 法とそのピボット版 CLIQUES、疎グラフ向けの Eppstein アルゴリズムを説明する。
続いて実験に用いる十種類のグラフを一覧で示し、最後に第\ref{ch:implementation}章以降で使う Python の multiprocessing の前提だけを述べる。

\section{グラフとクリーク}
\label{sec:graph-clique}

本論文では、無向グラフ $G = (V, E)$ を扱う。
$V$ は頂点の集合、$E \subseteq \{\, \{u, v\} \mid u, v \in V,\ u \neq v \,\}$ は辺の集合であり、頂点数を $n = |V|$、辺数を $m = |E|$ と書く。
自己ループと多重辺は持たないものとする。
頂点 $v$ の近傍を $N(v) = \{\, u \in V \mid \{u, v\} \in E \,\}$ で表し、その大きさ $|N(v)|$ を $v$ の次数と呼ぶ。

頂点部分集合 $C \subseteq V$ が\textbf{クリーク}であるとは、$C$ の任意の二頂点が辺で結ばれていること、すなわち $C$ の誘導部分グラフが完全グラフであることをいう。
クリーク $C$ が\textbf{極大クリーク}であるとは、$C$ にどの頂点を加えてもクリークでなくなること、言い換えれば $C \subsetneq C'$ となるクリーク $C'$ が存在しないことをいう。
極大クリークは、頂点数が最大のクリークを指す\textbf{最大クリーク}とは異なる概念である。
最大クリークは常に極大クリークだが、その逆は成り立たない。
本論文が対象とするのは、グラフ中のすべての極大クリークを重複なく出力する\textbf{極大クリーク列挙}である。

\section{Bron--Kerbosch 法とピボット版 CLIQUES}
\label{sec:bron-kerbosch}

極大クリーク列挙の基礎となるのが Bron--Kerbosch 法である。
この方法は、三つの頂点集合を持って再帰的に探索する。
確定済みの現在のクリーク $R$、まだ $R$ に加えられる候補 $P$、すでに調べ終えて再訪を禁じる $X$ である。
$P$ と $X$ がともに空になったとき、$R$ はそれ以上広げられず他の枝でも数えられない極大クリークなので、これを出力する。
候補 $P$ の各頂点 $v$ について、$v$ を $R$ に加え、$P$ と $X$ を $v$ の近傍だけに絞って再帰し、戻ったら $v$ を $P$ から $X$ へ移す。

素朴な Bron--Kerbosch 法は、同じ極大クリークに至る枝を何度も展開してしまう。
この無駄を抑えるのが\textbf{ピボット}の選択である。
候補と処理済みの合併 $P \cup X$ から一頂点 $u$（ピボット）を選ぶと、任意の極大クリークは $u$ かその非近傍のいずれかを必ず含むため、再帰を $P$ 全体ではなく「$u$ の非近傍」だけに限定できる。
Tomita, Tanaka, Takahashi は、$|P \cap N(u)|$ を最大にする $u$ をピボットに選ぶ版を\textbf{CLIQUES}として定式化した\cite{tomita2006}。
この選び方は再帰の分岐を最も強く刈り取り、最悪計算量を $O(3^{n/3})$ に抑える。

この計算量は、出力しうる極大クリークの個数という下限に照らして最善である。
$n$ 頂点の無向グラフが持ちうる極大クリークの個数は高々 $3^{n/3}$ であり、この上限は頂点を3個ずつの組に分けた完全多部グラフ $K_{3,3,\dots,3}$ が達成することを Moon と Moser が示した\cite{moon1965}。
極大クリークを一つ出力するのに最低でも定数時間かかる以上、列挙にかかる時間はこの個数を下回れない。
CLIQUES の $O(3^{n/3})$ はこの上限と一致しており、最悪の入力に対して漸近的に最適である\footnote{全体計算量と、極大クリークを一つ出すごとの遅延（delay）の関係は Conte と Tomita による後続の解析で詳しく調べられている\cite{conte2022}。本論文の逐次実装はこの論文の擬似コードに従う。}。

\section{縮退数と縮退順序}
\label{sec:degeneracy}

CLIQUES の最悪計算量は頂点数 $n$ だけで決まり、グラフが疎であることを利用しない。
現実の多くのグラフは辺が少なく、この疎さを計算量に反映できれば列挙は速くなる。
その疎さを測る尺度が\textbf{縮退数}である。

グラフ $G$ が $d$-縮退であるとは、$G$ のどの部分グラフにも次数 $d$ 以下の頂点が存在することをいう。
これが成り立つ最小の $d$ を $G$ の縮退数と呼ぶ。
縮退数は\textbf{縮退順序}を通しても特徴づけられる。
次数が最小の頂点を繰り返し取り除き、取り除いた順に並べたものが縮退順序であり、この順序では、どの頂点も自分より後ろに位置する近傍（後方近傍）を高々 $d$ 個しか持たない。
縮退順序は最小次数の頂点を順に剥がすだけで $O(n + m)$ 時間で構成でき、疎なグラフほど $d$ が小さい。

\section{Eppstein アルゴリズム}
\label{sec:eppstein}

Eppstein, L\"offler, Strash は、縮退順序を使って CLIQUES の計算量を縮退数で表す変種を与えた\cite{eppstein2010}。
最外層の再帰を縮退順序に沿って次数の低い頂点から回し、各頂点ではその後方近傍だけを候補に取る。
後方近傍は高々 $d$ 個なので、各頂点を根とする部分問題は最大 $d$ 頂点のグラフ上の列挙に帰着し、その内部ではピボット選択を Tomita 流のまま用いる。
これにより全体の計算量は $O(d \cdot n \cdot 3^{d/3})$ となり、$d$ が小さい疎グラフでは $n$ にほぼ比例する時間まで下がる。
実データを含む大規模な疎グラフでの有効性は、続く実験版の論文で確かめられている\cite{eppstein2013}。

このアルゴリズムの構造は、本論文の並列化にとって重要である。
最外層が縮退順序の各頂点を根とする $n$ 個の独立した部分問題に分かれるため、部分問題を worker へ配るだけで最外層を並列化できる。
$d$ が小さいほど各部分問題は小さく独立性が高い。

\section{実験に用いるグラフ}
\label{sec:graphs}

実験では、次数分布と疎密の異なる十種類の無向グラフを用いる。
生成規則が既知の合成グラフ七種と、実データである SNAP データセット三種に分かれる。
比較の主軸は縮退数 $d$ であり、$d$ が小さいグラフほど疎で、Eppstein アルゴリズムが最外層を小さく独立した部分問題に分割しやすい。

合成グラフ七種は、疎密に加えて次数分布と出力サイズという二つの軸でも選んでいる。
次数分布については、次数がほぼ均一な Erd\H{o}s--R\'enyi と、少数の頂点に次数が集中する Barab\'asi--Albert（ソーシャルネットワークに近いべき乗則分布）を対比させる。
出力サイズについては、辺の数だけ極大クリークが出る2次元格子と、理論上限 $3^{k}$ 個をそのまま達成する Moon--Moser を両端に置く。
各グラフの内容は次のとおりである。

\begin{itemize}
  \item \textbf{木}：頂点 $i$ が、それより前の頂点から一様ランダムに一つ選んで接続する再帰的ランダム木。三角形を持たず、$d = 1$。
  \item \textbf{2次元格子}：$100 \times 100$ の4近傍格子。三角形を持たないため、辺の一本ずつがそのまま極大クリークになる。$d = 2$。
  \item \textbf{疎 Erd\H{o}s--R\'enyi}：平均次数を10に固定して頂点対を一様ランダムに結ぶ $G(n, m)$ 型のグラフ。$d = 7$。
  \item \textbf{スケールフリー Barab\'asi--Albert}：新規頂点が既存頂点の次数に比例した確率で $m = 5$ 本の辺を張る優先的接続モデル。$d = 5$。
  \item \textbf{密 Erd\H{o}s--R\'enyi}：頂点対ごとに確率 $p = 0.5$ で辺を張る $G(n, p)$ モデル。辺数が $n^2$ で効くため小さい $n$ でのみ使う。$d = 48$。
  \item \textbf{Moon--Moser}：頂点を3個ずつ $k = 11$ 組に振り分けた完全多部グラフ $K_{3,3,\dots,3}$。頂点数 $3k$ に対し極大クリークが $3^{k}$ 個生じる、\ref{sec:bron-kerbosch}節で述べた極値グラフである。$d = 30$。
  \item \textbf{完全グラフ}：全頂点対に辺がある $K_{300}$。極大クリークは全体で1個だけになる。$d = 299$。
\end{itemize}

実データは、Stanford の SNAP（Stanford Large Network Dataset Collection）が配布する無向グラフ三種である。
いずれも頂点対の行を並べたエッジリスト形式で読み込み、有向や符号つきの情報は捨てて無向グラフとして扱う。
ca-AstroPh と ca-HepPh は arXiv の共著関係のネットワークで、それぞれ天体物理学分野と高エネルギー物理学分野に対応する。
soc-sign-epinions は商品レビューサイト Epinions のユーザー間の信頼・不信ネットワークで、符号を無視して無向グラフとして扱う。
このなかで soc-sign-epinions は頂点数と極大クリーク数がともに大きく、逐次で約95秒かかるため、第\ref{ch:experiments}章では並列化の効果を測る主力データとして用いる。

各グラフの頂点数 $n$、縮退数 $d$、極大クリーク数を表\ref{tab:graphs}に示す。

\begin{table}[htbp]
  \centering
  \caption{実験に用いるグラフ。上段が合成グラフ七種、下段が SNAP データセット三種。}
  \label{tab:graphs}
  \begin{tabular}{lrrr}
    \toprule
    グラフ & $n$ & $d$ & 極大クリーク数 \\
    \midrule
    木 & 20{,}000 & 1 & 20.0千 \\
    2次元格子 $100 \times 100$ & 10{,}000 & 2 & 19.8千 \\
    疎 Erd\H{o}s--R\'enyi（平均次数10） & 20{,}000 & 7 & 99.7千 \\
    スケールフリー Barab\'asi--Albert（$m=5$） & 20{,}000 & 5 & 96.9千 \\
    密 Erd\H{o}s--R\'enyi（$n=120$, $p=0.5$） & 120 & 48 & 3.5万 \\
    Moon--Moser（$k=11$） & 33 & 30 & 17.7万 \\
    完全グラフ $K_{300}$ & 300 & 299 & 1 \\
    \midrule
    ca-AstroPh & 18{,}772 & 56 & 36.4千 \\
    ca-HepPh & 12{,}008 & 238 & 14.9千 \\
    soc-sign-epinions & 131{,}828 & 121 & 2{,}223万 \\
    \bottomrule
  \end{tabular}
\end{table}

\todo{合成グラフ七種の外観図（Notion に描画済み画像あり）を挿入する。}

\section{Python multiprocessing による並列化}
\label{sec:multiprocessing}

第\ref{ch:implementation}章以降の並列化は Python の標準ライブラリ multiprocessing を用いる。
CPython はインタプリタ全体を一つのロック（GIL）で保護しており、複数スレッドが Python バイトコードを同時に実行できない。
そのため CPU バウンドな列挙を並列化するには、スレッドではなく複数のプロセスを使い、各プロセスに独立したインタプリタを持たせる必要がある。
以下では、並列化のコストを左右するプロセス生成方式の違いだけを述べる。

子プロセスの生成方式には \texttt{spawn} と \texttt{fork} がある。
\texttt{spawn} は子プロセスで新しいインタプリタを起動し、worker に渡す引数を親側で pickle により直列化して送り、子側で復元する。
\texttt{fork} は親プロセスのメモリを copy-on-write で子に継承させるため、引数の直列化も転送も起こらない。
macOS の既定は \texttt{spawn} であり、これは worker へ渡すグラフを worker 数だけ pickle 転送することを意味する。
このグラフ転送が並列化の固定費として効いてくる点は、第\ref{ch:experiments}章で実測とともに扱う。

\chapter{関連研究}
\label{ch:related-work}

\section{逐次アルゴリズムの系譜}
\label{sec:related-sequential}

極大クリーク列挙の逐次アルゴリズムは、Bron--Kerbosch 法を起点とする一つの系譜をなす。
その技術的な内容は第\ref{ch:preliminaries}章で述べたとおりである。

Tomita, Tanaka, Takahashi は、ピボット選択によって重複した枝の展開を防ぐ CLIQUES を提案し、その最悪計算量が極大クリーク数の理論上限 $3^{n/3}$ と一致することを示した\cite{tomita2006}。
この上限自体は、Moon と Moser が示した極値グラフによって与えられている\cite{moon1965}。
Conte と Tomita は、CLIQUES と Bron--Kerbosch 法の各変種について、全体計算量だけでなく極大クリークを一つ出力するごとの遅延（delay）まで踏み込んで解析し直した\cite{conte2022}。
本論文の CLIQUES 実装は、この論文が与える擬似コードに従っている。

Eppstein, L\"offler, Strash は、CLIQUES のピボット選択自体は保ちながら、最外層だけを縮退順序に沿って処理することで、計算量を縮退数 $d$ を用いた $O(d n 3^{d/3})$ に抑える変種を示した\cite{eppstein2010}。
Eppstein と Strash は、続く実験で、実データを含む大規模な疎グラフでもこの変種が CLIQUES より高速に動作することを確かめている\cite{eppstein2013}。
本論文は、この Eppstein のアルゴリズムを、CLIQUES との逐次比較を経た上で並列化の対象に選んでいる。
選定の根拠となる測定結果は第\ref{ch:experiments}章で示す。

\section{並列化に関する先行研究}
\label{sec:related-parallel}

MCE の並列化についても複数の先行研究がある。

Schmidt, Hadfield, Samatova, Rice, Schuchardt は、共有メモリ環境で動作するスケーラブルな並列 MCE アルゴリズムを提案した\cite{schmidt2009}。
Das, Svendsen, Tirthapura による\textbf{ParMCE} は、共有メモリ上で動作する並列 MCE アルゴリズムであり、32 コア環境で逐次実行比最大 31 倍の高速化を報告している\cite{das2018parmce}。
グラフ分割に基づく別の手法も、既存手法比で最大 10 倍の高速化を報告している\cite{jpdc2017kplex}。
これらの論文はいずれも抄録レベルでしか内容を確認できておらず、具体的な実装は本文を未読のままである。
\todo{各論文の本文を確認し、アルゴリズムと実験条件の詳細を加筆する}

抄録から読み取れる範囲でも、これらの先行研究には共通する設計上の定石が見て取れる。
第一に、探索木をあらかじめプロセッサ数よりずっと多いサブツリーに分割しておく\textbf{オーバーデコンポジション}である。
第二に、暇になった worker が未処理のサブツリーを実行時に取りに行く\textbf{work stealing}（動的スケジューリング）であり、負荷の偏りを実行しながら解消する。
第三に、グラフ本体を worker ごとにコピー転送するのではなく共有メモリ上に置く配置であり、コピーにかかる固定費そのものを発生させない。

\section{本研究の位置づけ}
\label{sec:related-position}

本論文の並列実装は、これらの定石とは異なり、縮退順序で並べた頂点をバッチに区切って worker へ静的に配る素朴な配分から出発している\footnote{バッチの並べ方（block、reversed、interleave の3種）は第\ref{ch:implementation}章で述べる。}。
静的な配分は、work stealing のような実行時の負荷再分配を持たないため、負荷の偏りをバッチの作り方でしか吸収できない。
それでも、第\ref{ch:experiments}章で示すとおり、バッチの並べ方を変えるだけで worker の遊休はほぼ解消できた。
一方で、遊休を解消しても実行時間の短縮は先行研究が報告するような大幅な高速化には届かず、その理由は性能コアと効率コアの非対称な構成や共有資源の競合という、配分の工夫だけでは解決できない要因にあった。

先行研究は、これらの定石がなぜ有効かという機構までは抄録から読み取れない。
本論文は、単一マシン上での実測によって、work stealing とオーバーデコンポジションが解消しようとしている負荷の偏りの実体（重い頂点が縮退順序の末尾に密集すること）と、共有メモリ配置が解消しようとしている固定費の実体（プロセス起動時のグラフ転送コスト）を、それぞれ切り分けて確認する。
新しい並列アルゴリズムを提案するのではなく、既存の定石が効く理由とその限界を単一マシンの実測によって裏づける、再現・分析型の研究として本論文を位置づける。

\chapter{並列化の設計と実装}
\label{ch:implementation}

本章では、第2章で述べた Eppstein アルゴリズムを並列化するための設計と実装を述べる。
実装はどれも同じ骨格を共有し、変える点を一つに絞った変種の集まりとして構成した。
バッチの配り方だけを変えた三つの変種、グラフを worker へ届ける方式だけを変えた二つの変種であり、性能差の原因を一つの設計判断に帰着できるようにするためである。
各変種の実測結果は第5章で示す。

\section{最外層の独立性とカウント版への限定}

並列化の足場は、Eppstein アルゴリズムの最外層が持つ独立性にある。
このアルゴリズムは縮退順序に沿って各頂点 $v_i$ を一度ずつ根に取り、その頂点を含む極大クリークだけを担当する部分問題を解く。
頂点 $v_i$ の部分問題は、固定されたグラフと縮退順序だけに依存し、ほかの頂点の部分問題の結果には一切依存しない。
したがって最外層のループは、そのまま互いに独立なタスクの集まりとして切り出せる。

この独立性を、それ以上の工夫を足さずにそのまま並列化の単位とした。
再帰の内側（ピボット選択と候補集合の絞り込み）には手を入れず、外側の反復だけを複数のプロセスに振り分ける。
内側にも並列性はあるが、部分問題どうしがすでに独立している以上、まず最外層だけを配る素朴な構成で固定費と負荷分布を測り、そこで見えた限界に応じて次を決める方針を採った。

並列化の対象は、極大クリークの個数を数える\textbf{カウント版}に限定した。
列挙版は極大クリークの集合そのものを出力するが、soc-sign-epinions では出力が2千万個を超える。
これを複数プロセスから集約するには、各 worker が見つけたクリークをプロセス間で転送してまとめる必要があり、列挙結果の転送と集約という別の重い問題が主役になってしまう。
本研究が測りたいのは列挙計算そのものの並列化なので、出力は各 worker がローカルに数えた整数の総和で済むカウント版を対象とし、転送コストを個数の足し合わせだけに切り離した。

\section{プロセスプールによる基本実装}

基本実装は、Python 標準ライブラリの \verb|multiprocessing.Pool| による素朴なプロセスプールである。
縮退順序の計算は親プロセスで逐次に一度だけ行い、その順序を連続する \verb|BATCH_SIZE| 頂点ずつのバッチに区切って worker へ配る。
各 worker は担当バッチ内の頂点それぞれについて、近傍のうち縮退順序で後ろにあるものを候補集合、前にあるものを除外集合として、ピボット再帰で部分木の極大クリーク数を数える。

worker へのデータ配布には、プールの初期化関数を使う。
\verb|Pool| の \verb|initializer| でグラフと縮退順序を各 worker のモジュール大域変数に一度だけ格納し、以降のタスクはバッチの範囲だけを受け取る。
こうしないと、グラフをタスクごとに pickle 化して送り直すことになり、転送コストがタスク数だけ積み上がる。
macOS のプロセス起動方式は既定で spawn なので、この初期データはプール起動時に pickle 化されてパイプ経由で各 worker へ転送される。
グラフ全体を worker ごとに一度コピーするこの設計は、後で固定費として効いてくる（第5章）が、まずは最も素朴な形を基準線として測るために選んだ。

負荷分散は \verb|imap_unordered| の動的配分に任せた。
バッチは完成した順に結果を回収でき、速く終わった worker から次のバッチを引く。
静的に「どの worker がどのバッチを担当するか」を決め打ちせず、実行時のタスク取得だけで均す狙いである。

\begin{verbatim}
with Pool(workers, initializer=_init_worker,
          initargs=(graph, ordering)) as pool:
    pool.map(_noop, range(workers * 4), chunksize=1)  # 起動バリア
    for pid, elapsed, sub_count, largest in pool.imap_unordered(
            _count_batch_timed, batches):
        count += sub_count
\end{verbatim}

\verb|BATCH_SIZE| は、動的配分が働くだけの細かさと、タスク発行のオーバーヘッドが無視できるだけの粗さの折り合いで決まる。
バッチが大きすぎると、重いバッチを引いた worker が一台だけ長く走り、ほかが待つ。
逆に細かすぎれば、タスクの受け渡し回数が増える。
この粒度が並列化の効きを左右することは第5章で示す。

\section{バッチ配分の三つの戦略}

基本実装のバッチは縮退順序の先頭から連続に区切るが、この区切り方には偏りの問題がある。
負荷分布の測定（第5章）で、重い頂点は縮退順序の末尾に密集することがわかっている。
疎な外周から頂点を剥がしていく縮退順序では、密なコアが必然的に末尾へ回るためである。
連続区切りだと重い頂点が同じ末尾のバッチにまとまり、そのバッチを引いた worker だけが突出して長く走る。

この偏りにバッチの作り方だけで対処するため、区切り方の異なる三つの戦略を用意した。
枠組みはどれも「縮退順序の各頂点を根とする部分問題を、バッチにまとめて worker に配る」で共通し、違うのはバッチの中身の決め方だけである。

\begin{itemize}
\item \textbf{block}（基本実装）：順序の先頭から連続に区切る。重いバッチが末尾にまとまり、最後に配られる。
\item \textbf{reversed}（末尾逆順）：順序を逆にしてから連続に区切る。重いバッチを最初に配り、終盤で重いバッチを引いて長く待つ事態を防ぐ。
\item \textbf{interleave}（ストライド）：バッチ数を $K$ として、バッチ $k$ に順序上の頂点 $k, k+K, k+2K, \dots$ を入れる。どのバッチにも軽い頂点と重い頂点が混ざり、バッチの重さ自体がほぼ均一になる。
\end{itemize}

\begin{verbatim}
# block:      batches = [range(i, min(i + b, n)) for i in range(0, n, b)]
# reversed:   idx = list(range(n))[::-1]
#             batches = [idx[i:i+b] for i in range(0, n, b)]
# interleave: K = -(-n // b)  # バッチ数
#             batches = [range(k, n, K) for k in range(K)]
\end{verbatim}

reversed と interleave は、同じ偏りに別の角度から当たっている。
reversed は重い順に配って終盤の待ちを消すが、最重のバッチ1個より速くは終われない。
interleave は各バッチの重さをならして、最重のバッチ自体を軽くする。
どちらが効くか、そもそも静的な配り方で偏りを救えるのかは第5章で検証する。

\section{グラフを worker へ届ける二つの方式}

基本実装の固定費は、グラフを spawn で各 worker へ pickle 転送するところに集中する。
この固定費を削るために、グラフの届け方だけを変えた二つの変種を用意した。
どちらもバッチの作り方は block と同一にし、block との差がすべて配布方式に帰着するようにしてある。

\subsection*{fork による copy-on-write 継承}

一つ目は、プロセス起動方式を spawn から fork に変える変種である。
\verb|multiprocessing.get_context("fork")| で worker を起動すると、子プロセスは親のメモリを copy-on-write で継承する。
グラフは親のアドレス空間にすでにあるので、pickle 化もパイプ転送も worker 側の再構築も起きず、起動コストは fork 自体まで縮むはずである。

macOS が既定で fork を使わないのは、スレッドを持つプロセスを fork すると、ロックを保持したままのスレッドが子には複製されず、子がデッドロックしうるためである。
本実装は fork を呼ぶ時点で単一スレッドの純 Python であり、この危険がないので fork を安全に使える。
既定を外す判断の根拠がこの単一スレッド性にあることを、変種の docstring にも明記した。

\subsection*{shared\_memory による共有配列}

二つ目は、起動方式は spawn のまま、グラフを共有メモリに置く変種である。
親プロセスがグラフを CSR 形式の int64 配列4本（頂点 ID、隣接リストの開始位置、隣接先、縮退順序）にまとめ、\verb|multiprocessing.shared_memory| のブロックに書き込む。
各 worker は転送を受け取る代わりに、ブロックへ名前で attach する。

ただし、この変種でも worker 側の構造構築は消えない。
再帰カーネルは Python の集合演算（\verb|set|）に依存し、\verb|set| は共有メモリに置けないためである。
共有メモリに置けるのは連続したバイト列だけで、ハッシュテーブルである \verb|set| の内部構造は載らない。
そのため各 worker は、attach した共有配列から自分用の集合の辞書を起動時に組み立て直す。

\begin{verbatim}
graph = {ids[i]: {ids[j] for j in indices[indptr[i]:indptr[i + 1]]}
         for i in range(len(ids))}   # 共有 CSR 配列から set を再構築
\end{verbatim}

この設計により、shm 変種が消すのは pickle 化とパイプ転送だけになる。
spawn の起動コストを「インタプリタ起動」「転送」「構造構築」の三つに分けたとき、転送だけを消すのが shm、三つとも消すのが fork という対比になる。
shm を単独の高速化策としてではなく、spawn の固定費の内訳を切り分ける測定器として位置づけたのはこのためである。

\section{フェーズ分解と worker 稼働時間の計測}

並列化の効き方を診断するために、1回の実行を四つの区間に分けて壁時計時間を測る計測器を実装した。
区間は、グラフの読込または生成（load）、縮退順序の計算（ordering、逐次）、プールの生成と worker 初期化（startup）、並列本体（compute）である。
load と ordering は逐次項なので Amdahl の法則における並列化の上限に効き、startup は worker 数に応じて伸びる固定費、compute は実際に並列で走る部分にあたる。
この四分割により、実行時間がどの区間に消えているかを worker 数ごとに追える。

startup の測り方には近似が含まれる。
\verb|Pool()| は子プロセスがグラフの unpickle を終える前に制御を返すため、プール生成の呼び出しが返った時刻は初期化の完了時刻ではない。
そこで、全 worker へ \verb|chunksize=1| で noop タスクをまき、それらが返るまでを起動バリアとした。
これが近似にとどまるのは、\verb|imap_unordered| が後続のバッチを、noop のまき方によっては初期化直後の特定の worker に偏って渡しうるためで、厳密な転送完了時刻とは一致しない。
この近似の性格は、startup の絶対値を手法間で比較するときの限界として第5章でも断る。

compute 区間では、worker ごとの稼働時間も同じ実行から取得する。
各バッチ処理が自プロセスの PID と所要時間を返し、親が PID ごとに積算する。
これにより、最も忙しい worker と最も暇な worker の稼働時間の差、および遊休率を測れる。
遊休率は、全 worker の稼働時間の合計を「worker 数 $\times$ compute 時間」で割った値を1から引いて求める。
バッチを一つも受け取らなかった worker も処理能力としては存在するので、分母は稼働した worker 数ではなく起動した worker 数を使う。
この遊休率が、負荷分布の偏りが worker の遊びとして現れる度合いを示す指標になる。

\todo{第5章（実験）で示す測定結果への具体的な相互参照（節番号・図表番号）を、実験章の確定後に補う}

\chapter{実験と考察}
\label{ch:experiments}
\label{chap:experiments}

本章では、Eppstein アルゴリズムの最外層並列化を実装し、その性能を測定した過程をたどる。
話の順序は、実際に測定を進めた時系列にそろえてある。
ひとつの測定が次の問いを生み、その問いに答えるために次の測定を設計する、という連鎖として書く。
逐次の二アルゴリズムの比較から始め、素朴な並列化がなぜ頭打ちになるかを負荷分布とフェーズ別の内訳で突き止め、配分戦略・非対称コア・共有メモリ化と順に切り込んでいく。

各節は、何を知りたいか（意図）、どう測ったか（方法）、何が出たか（結果）、それが何を意味するか（考察）の順で述べる。
測定値には、直接測ったものと、第三者の測定や単純なモデルから導いた推定とがある。
両者は文中で区別し、推定にはその根拠を添える。

\section{実験環境と測定プロトコル}
\label{sec:exp-env}

\subsection{意図}

並列性能の測定は環境に敏感である。
本研究のように単一のノート型計算機で測る場合、背景で動くアプリケーションや熱の状態が数値を大きく揺らす。
どの数値がアルゴリズムの性質を表し、どの数値が測定環境の副作用かを後から切り分けられるよう、測定条件と手続きをここで固定しておく。

\subsection{実験環境}

測定はすべて Apple M4 を搭載した MacBook 1 台で行った。
M4 は高速な Performance コア（以下 P コア）4 個と低速な Efficiency コア（以下 E コア）6 個の非対称構成で、主記憶は 16\,GB である。
実装は外部の数値ライブラリを用いない純 Python で、並列化には標準ライブラリの \texttt{multiprocessing} を使う。
P コアと E コアの速度差、および 16\,GB という主記憶の制約は、後の節（\ref{sec:pe-split}、\ref{sec:ecore-speed}）で性能の説明そのものに関わってくる。

\subsection{測定プロトコル}

計測値の揺れに対しては、次の三つの手続きで対処する。

\begin{itemize}
  \item \textbf{反復最小値}：同一条件を 3 回繰り返し、最小値を採用する。実行時間は他プロセスの割り込みで上振れする一方、下振れはしない。最小値は割り込みの影響が最も小さい実行に対応する。
  \item \textbf{環境スナップショット}：測定の直前に \texttt{ps} で背景プロセスの状態を記録し、想定外の負荷が走っていないことを確認してから実行する。
  \item \textbf{時間的な隣接}：比較したい条件どうしは、同一セッション内で時間を空けずに連続して測る。後述する熱ドリフトが条件間に混入するのを防ぐためである。
\end{itemize}

このプロトコルが必要なことは、環境の影響を定量した二つの観測が裏づける。
ひとつは背景アプリの影響で、Chrome を閉じただけで soc-sign-epinions の並列列挙時間が約 26\% 短縮した（\ref{sec:pe-split} 節で詳述）。
もうひとつは熱ドリフトで、約 1.5 時間の連続測定のあと、同一の純計算ループが 0.64 秒から 0.85 秒へ 32\% 遅くなった（\ref{sec:ecore-speed} 節）。
どちらもアルゴリズムとは無関係な要因であり、この幅を条件間の差と取り違えないための手続きが上の三点である。

\section{逐次版の比較：Tomita 対 Eppstein}
\label{sec:seq-compare}

\subsection{意図}

並列化の土台として、まず逐次の実装を比較する。
本研究は疎グラフに有利な Eppstein を並列化の対象に選ぶが、その選択が妥当かを 10 種類のグラフで確かめておく。
比較の軸は\textbf{縮退数} $d$ である。
$d$ は、どの頂点も縮退順序で自分より後ろの近傍が $d$ 個以下になるように頂点を並べられる最小値で、グラフの疎さを表す。
Eppstein は最外層をこの縮退順序で回ることで計算量を $O(d\,n\,3^{d/3})$ に抑える。

\subsection{方法}

比較対象は、Tomita らのピボット版 CLIQUES と Eppstein の 2 実装である。
両者を同一プロセス内で交互に 3 回ずつ実行し、それぞれの最小値を採る。
交互反復にするのは、このマシンでは単発計測が 2〜3 割振れることがあるためで、交互に測れば環境の変動が両者に等しく乗り、相対比較が安定する。
3 回の振れは全データセットで 1\% 未満だった。
グラフは、生成規則が既知の合成グラフ 7 種と、Stanford の SNAP が配布する実データ 3 種からなる。

\subsection{結果}

結果を表~\ref{tab:seq-compare} に示す。
倍率は Tomita の時間を Eppstein の時間で割った値で、1 を超えれば Eppstein が速い。

\begin{table}[htbp]
  \centering
  \caption{逐次版の比較（各値は 3 回の最小値、倍率は Tomita/Eppstein）}
  \label{tab:seq-compare}
  \begin{tabular}{lrrrrrr}
    \toprule
    グラフ & $n$ & $d$ & 極大クリーク数 & Tomita & Eppstein & 倍率 \\
    \midrule
    木                       & 20{,}000  & 1   & 20.0千   & 0.502\,s & 0.030\,s & 16.6 \\
    2 次元格子 $100\times100$ & 10{,}000  & 2   & 19.8千   & 0.137\,s & 0.022\,s & 6.2 \\
    疎 Erdős–Rényi（平均次数10）& 20{,}000 & 7   & 99.7千   & 0.621\,s & 0.142\,s & 4.4 \\
    スケールフリー BA（$m{=}5$） & 20{,}000 & 5   & 96.9千   & 0.717\,s & 0.135\,s & 5.3 \\
    ca-AstroPh               & 18{,}772  & 56  & 36.4千   & 2.290\,s & 0.871\,s & 2.6 \\
    ca-HepPh                 & 12{,}008  & 238 & 14.9千   & 1.008\,s & 0.636\,s & 1.6 \\
    soc-sign-epinions        & 131{,}828 & 121 & 2{,}223万 & 167.3\,s & 94.9\,s & 1.8 \\
    密 Erdős–Rényi（$n{=}120$）& 120      & 48  & 3.5万    & 0.123\,s & 0.109\,s & 1.1 \\
    Moon–Moser（$k{=}11$）    & 33        & 30  & 17.7万   & 0.147\,s & 0.120\,s & 1.2 \\
    完全グラフ $K_{300}$      & 300       & 299 & 1       & 0.074\,s & 0.192\,s & 0.4 \\
    \bottomrule
  \end{tabular}
\end{table}

\subsection{考察}

結果は $d/n$（縮退数がグラフサイズに占める割合）でほぼ説明できる。
疎で大きいグラフほど Eppstein が勝ち、密になると互角に収束し、明確に負けるのは完全グラフだけである。
木（$d{=}1$）では 16.6 倍、完全グラフ（$d{=}299$）では 0.4 倍と、勝敗が $d$ の小ささにそのまま従っている。
Eppstein が最外層を縮退順序で回って最外層を小さな部分問題に分割できるのは、$d$ が小さい疎グラフに限られ、密グラフでは分割の余地がないためこの利点が消える。

以降の並列化はすべて Eppstein を対象とする。
主力のベンチマークには soc-sign-epinions を選ぶ。
逐次で約 94.8 秒かかり、並列化の効果を測るのに十分重く、かつ実データである点が理由である。

\section{素朴な並列化とバッチサイズ}
\label{sec:naive}

\subsection{意図}

最外層の各頂点を根とするサブ問題は互いに独立に処理できる。
そこで最外層で素朴に並列化し、worker 数を増やしたときにどこまで速くなるかを測る。

\subsection{方法}

グラフはプール起動時に各 worker へ一度だけ \texttt{pickle} 転送し、負荷分散は \texttt{imap\_unordered} の動的配分に任せる。
サブ問題はバッチにまとめて配る。
まずバッチサイズ 512 で worker 数を 1, 2, 4, 8 と変えて測った。

\subsection{結果}

結果を表~\ref{tab:naive-batch} に示す。

\begin{table}[htbp]
  \centering
  \caption{soc-sign-epinions の素朴な並列化（列挙時間、workers=1 比）}
  \label{tab:naive-batch}
  \begin{tabular}{lrr}
    \toprule
    構成 & 列挙時間 & workers=1 比 \\
    \midrule
    workers=1（基準）        & 99.8\,s & 1.00 \\
    workers=2、バッチ512     & 69.3\,s & 1.44 \\
    workers=4、バッチ512     & 67.1\,s & 1.49 \\
    workers=8、バッチ512     & 70.0\,s & 1.43 \\
    workers=8、バッチ64      & 27.7\,s & 3.60 \\
    \bottomrule
  \end{tabular}
\end{table}

バッチ 512 では 4 worker 以降がまったく伸びず、67〜70 秒で頭打ちになった。
変更したのはバッチサイズを 512 から 64 に細かくしたことだけで、それにより 8 worker が 70 秒から 27.7 秒まで縮んだ。

並列化が逆効果になる場合も確認した。
1 秒未満で終わるグラフ（例：木 $n{=}20{,}000$）では、プール起動とグラフの \texttt{pickle} 転送（合計約 0.3 秒）が支配的になり、すべてのケースで逐次より遅くなった。
ca-HepPh は固定費に加えて、239 頂点のクリークを含む探索木が 1 タスクに載るため、8 worker でも改善しなかった。

バッチ 512 の頭打ちは、負荷の偏りを示唆する。
そこで、頭打ちの理由を負荷分布の測定で確かめる。

\section{負荷分布の測定}
\label{sec:load-dist}

\subsection{意図}

バッチ 512 で並列化が頭打ちになった理由を明らかにする。
外側頂点ごとのサブ問題の所要時間を、全頂点で個別に計測した。

\subsection{結果}

結果を表~\ref{tab:load-dist} に示す。
数値は全所要時間に占めるシェアである。

\begin{table}[htbp]
  \centering
  \caption{サブ問題の負荷分布（全所要時間に占めるシェア）}
  \label{tab:load-dist}
  \begin{tabular}{lrrrrr}
    \toprule
    データセット & 頂点数 & 最重の1頂点 & 上位0.1\% & 上位1\% & 上位10\% \\
    \midrule
    soc-sign-epinions & 131{,}828 & 0.86\% & 45.2\% & 94.7\% & 99.7\% \\
    ca-HepPh          & 12{,}008  & 13.0\% & 20.0\% & 47.0\% & 90.1\% \\
    ca-AstroPh        & 18{,}772  & 0.21\% & 2.8\%  & 17.7\% & 65.9\% \\
    \bottomrule
  \end{tabular}
\end{table}

\subsection{考察}

soc では上位 1\% の頂点が仕事の 94.7\% を占め、その重い頂点は縮退順序の末尾に密集していた。
疎な外周から頂点を剥がしていくと密なコアが必然的に末尾に残るためである。
ただし最重の 1 頂点でも 0.86\% なので、バッチを細かくすれば負荷は均せる。
バッチ 512 から 64 への変更で 8 worker が 70 秒から 27.7 秒に縮んだのは、この偏りの解消にあたる。

ca-HepPh は性質が異なる。
1 頂点（239 頂点クリークの根）だけで全時間の 13.0\% を占め、この 1 タスクは分割しない限りどう配っても待たされる。
このためスピードアップの上限が約 7.7 倍に制限される。
「配り方で救える soc」と「サブ問題内部の分割が要る ca-HepPh」は、同じ「並列化が効かない」でも別の問題である。
バッチの細分化で救えるのは前者だけで、後者はタスクの粒度そのものが下限を作る。

\section{フェーズ別時間の分解}
\label{sec:phase}

\subsection{意図}

並列実行の時間がどこに消えているかを知る。
1 回の実行を 4 区間に分けて計測した。

\begin{itemize}
  \item \textbf{load}：グラフの読込・生成。並列実行の外側で 1 回だけ測る、worker 数に依存しない値。
  \item \textbf{ordering}：縮退順序の計算。逐次に行う準備ループ。
  \item \textbf{startup}：プロセスプールの生成と worker の初期化。
  \item \textbf{compute}：並列本体。
\end{itemize}

あわせて、同じ実行から worker 別の busy 時間と idle 率を取得した。
計測は 9 グラフ $\times$ workers=1,2,4,8 $\times$ 3 回（最小値採用）である。

\subsection{準備ループは軽い}

最初の問いは、逐次の準備ループ（縮退順序）がどれほど重いか、である。
答えは、心配するほど重くない、だった。
ordering が total に占める割合は全グラフ・全 worker 数を通して最大 17.8\%（絶対値 0.05 秒未満）で、主力の soc では 1\% 未満である。
Amdahl の逐次項としては当面無視してよい。

\subsection{固定費は startup に集中する}

かわりに固定費として効いていたのは startup だった。
プール起動時にグラフ全体を worker ごとに \texttt{pickle} 転送する設計の代償で、worker 数にほぼ比例して伸びる。
soc では w=1 の 0.27 秒から w=8 の 1.27 秒まで増える。
compute が 1 秒未満の 7 グラフでは、この startup のために w=8 が w=1 より遅くなる。

soc の worker 数スイープをフェーズ別に分けた内訳を表~\ref{tab:phase-soc} に示す。

\begin{table}[htbp]
  \centering
  \caption{soc-sign-epinions のフェーズ別内訳（バッチ64、3 回の最小値）}
  \label{tab:phase-soc}
  \begin{tabular}{rrrrrrrl}
    \toprule
    workers & ordering & startup & compute & total & 対w=1倍率 & idle\% & busy 最大/最小 \\
    \midrule
    1 & 0.432\,s & 0.268\,s & 126.006\,s & 126.706\,s & 1.00 & 0.0\%  & 126.0 / 126.0\,s \\
    2 & 0.366\,s & 0.338\,s & 65.800\,s  & 66.504\,s  & 1.91 & 6.9\%  & 65.7 / 56.8\,s \\
    4 & 0.339\,s & 0.643\,s & 48.419\,s  & 49.401\,s  & 2.56 & 15.0\% & 48.4 / 35.6\,s \\
    8 & 0.338\,s & 1.273\,s & 36.193\,s  & 37.803\,s  & 3.35 & 30.4\% & 36.2 / 11.8\,s \\
    \bottomrule
  \end{tabular}
\end{table}

startup が worker 数にほぼ比例して伸びているのが直接読める。
startup の値は近似である。
プール生成の直後に全 worker へ noop を配って初期化完了を待つ方式で測っているため、厳密な転送完了時刻とは一致しない。

\subsection{並列化が得なグラフは限られる}

9 グラフを通して見ると、並列化で得をしたのは逐次 compute が長い 3 グラフだけだった。
soc は worker を増やすほど伸びる（w=8 で 3.35 倍）のに対し、ca-AstroPh と ca-HepPh は w=4 がピーク（1.62 倍、1.51 倍）で、w=8 では startup の伸びが compute の削減を上回って逆に悪化する。
残りの 6 グラフ（compute が 0.2 秒以下）はすべての worker 数で損で、最悪の 2 次元格子は w=8 で w=1 の 5 倍遅い。

\subsection{スケールしない主因は compute 内の idle}

soc に話を戻すと、w=8 でも ordering と startup の合計は total の 4\% に過ぎない。
スケールしない主因は固定費ではなく、compute 内部の idle である。
idle\% は w=2 の 6.9\% から w=8 の 30.4\% へ伸び、busy 時間は最も忙しい worker の 36.2 秒に対し最も暇な worker は 11.8 秒だった。
負荷分布の測定で見つけた「末尾の重い頂点への偏り」が、worker の遊休として現れている。

次の問いは、この compute 内の idle を配り方だけでどこまで消せるか、である。

\section{配分戦略の比較}
\label{sec:alloc}

\subsection{意図}

フェーズ別分解で見つけた compute 内の遊休（w=8 で idle 30.4\%、busy 最大 36.2 秒に対し最小 11.8 秒）を、タスクの配り方だけでどこまで解消できるかを検証する。
先行研究の定石である work stealing や共有メモリ化は実装コストが大きい。
その前に、静的な配分の工夫で届く範囲を確かめておく位置づけである。

\subsection{方法}

比較した三つの戦略は、いずれも「縮退順序の各頂点を根とするサブ問題を、バッチにまとめて worker に配る」枠組みは同じで、バッチの作り方だけが違う。
前提として、負荷分布の測定から重い頂点は縮退順序の末尾に密集することがわかっている。

\begin{itemize}
  \item \textbf{block}（現行）：順序の先頭から連続にバッチサイズずつ区切る。重いバッチが最後に配られる。
  \item \textbf{reversed}（末尾逆順）：順序を逆にしてから連続に区切る。重いバッチを最初に配り、終盤の待ちを防ぐ。
  \item \textbf{interleave}（ストライド）：バッチ数を $K$ として、バッチ $k$ に頂点 $k, k+K, k+2K, \dots$ を入れる。どのバッチにも軽い頂点と重い頂点が混ざり、バッチの重さ自体がほぼ均一になる。
\end{itemize}

測定条件は soc-sign-epinions、バッチ 64、workers=4,8、各条件 3 回の最小値である。
単発計測の振れ対策として、block も同一セッションで再測定した基準を使う。
極大クリーク数は全条件で 22{,}226{,}173 に一致し、結果の正しさは確認済みである。

\subsection{結果}

結果を表~\ref{tab:alloc} に示す。

\begin{table}[htbp]
  \centering
  \caption{配分戦略の比較（soc-sign-epinions、バッチ64、3 回の最小値）}
  \label{tab:alloc}
  \begin{tabular}{llrrrr}
    \toprule
    戦略 & workers & compute & idle\% & busy 最大/最小 & busy 合計 \\
    \midrule
    block      & 4 & 33.3\,s & 20.2\% & 33.3 / 22.2\,s & 106\,s \\
    reversed   & 4 & 32.7\,s & 0.1\%  & 32.7 / 32.7\,s & 131\,s \\
    interleave & 4 & 36.7\,s & 0.1\%  & 36.7 / 36.6\,s & 147\,s \\
    block      & 8 & 25.9\,s & 32.6\% & 25.8 / 8.0\,s  & 139\,s \\
    reversed   & 8 & 28.5\,s & 23.5\% & 28.5 / 18.1\,s & 175\,s \\
    interleave & 8 & 24.8\,s & 0.1\%  & 24.7 / 24.6\,s & 198\,s \\
    \bottomrule
  \end{tabular}
\end{table}

w=8 の最速は interleave の 24.8 秒で、block からの短縮は 4.3\% にとどまった。
reversed は w=4 では僅かに勝つが、w=8 では block より悪化した。

\subsection{負荷の均等化は成功し、高速化は失敗した}

interleave は idle を 32.6\% から 0.1\% までほぼ完全に消した。
idle が純粋な無駄なら、compute は $1/(1-0.326)$ で約 1.5 倍縮むはずだったが、実際の短縮は 4.3\% だった。
理由は表の busy 合計に表れている。
idle が消えるのと引き換えに、worker が働いた時間の合計が 139 秒から 198 秒へ膨張しており、「遊んでいた worker を働かせる」効果を「1 単位の仕事にかかる時間が伸びる」効果がほぼ打ち消している。

膨張の有力な説明は、測定マシンが P コア 4 個と E コア 6 個の非対称構成であることだ。
w=8 では必ず E コアにも worker が載るため、worker の稼働時間は仕事量の物差しにならない。
idle\% という診断指標は「コアは均質である」という前提を暗黙に置いており、その前提が崩れるこのマシンでは、負荷均等化の効果を過大に予測していた。
もうひとつの候補として、重いサブ問題どうしを同時に走らせることによるキャッシュとメモリ帯域の競合も busy 膨張に寄与しうるが、こちらの切り分けは未検証である。

reversed が w=8 で悪化した機構は別である。
最重バッチを最初に引いた worker の busy（28.5 秒）が compute 全体と一致しており、最重バッチ 1 個がクリティカルパスになっている。
重い方から配る戦略は、バッチをさらに細かくしない限りこの下限を破れない。

帰結として、静的配分の改善はこのマシンではほぼ頭打ちである。
idle をほぼ 0\% にしてもなお高速化しなかったという事実は、次の二つの主張を検証する動機になる。
第一に、E コア 6 個で走る分の処理が全体の足を引っ張っているのではないか。
第二に、主記憶 16\,GB が不足し、swap や競合が起きているのではないか。

\section{P/E コアの切り分け実験}
\label{sec:pe-split}

\subsection{意図}

前節の二つの主張を検証する。
主張 1（P/E 非対称コアが足を引っ張っている）は、macOS の \texttt{taskpolicy -b} でプロセスを E コアに固定して同じ計算を走らせれば、E コアの速度を直接測れる。
主張 2（16\,GB 不足による swap・競合）は、グラフのメモリフットプリントの実測で swap の有無を判定できる。
あわせて、Chrome を閉じた環境で interleave の w=4,8 を再測定し、背景アプリの影響も確かめる。

\subsection{方法}

条件は soc-sign-epinions、interleave、バッチ 64、各条件 3 回の最小値である。
すべて同一セッションで、Chrome を閉じ、測定前に \texttt{ps} で背景プロセスがアイドルであることを確認してから実行した。
E コア条件は \texttt{taskpolicy -b} 経由で起動する。
background QoS に落とすことで Apple Silicon では E コアに固定されるが、この方法は優先度も同時に下げるため、測った E コア速度は下限値として扱う。
逆に P コア側を強制する手段は macOS になく、「重い単一プロセスは P コアに載る」というスケジューラの性質に依存している。

\subsection{結果}

結果を表~\ref{tab:pe-split} に示す。
スループットは w=1 通常の compute（96.0 秒）を当該条件の compute で割った値で、「P コア何個分の仕事が出たか」を表す。
極大クリーク数は全条件で一致する。

\begin{table}[htbp]
  \centering
  \caption{P/E コアの切り分け（soc-sign-epinions、interleave、3 回の最小値）}
  \label{tab:pe-split}
  \begin{tabular}{lrrrr}
    \toprule
    条件 & compute & スループット（P換算） & worker 1個あたり & 単独比の効率 \\
    \midrule
    w=1 通常（P）           & 96.0\,s  & 1.00 & 1.00 & ― \\
    w=1 E コア固定          & 363.6\,s & 0.26 & 0.26 & ― \\
    w=4 通常（実質 P のみ）  & 30.4\,s  & 3.16 & 0.79 & 79\% \\
    w=4 E コア固定          & 124.5\,s & 0.77 & 0.19 & 73\% \\
    w=8 通常（P+E 混載）     & 18.2\,s  & 5.26 & 0.66 & ― \\
    \bottomrule
  \end{tabular}
\end{table}

\subsection{主張 1 は実証された}

同じ計算が P コアで 96 秒、E コア固定で 363.6 秒と約 4 倍違う（速度比 0.26、ただし background QoS の影響を含む下限値）。
worker の稼働時間が仕事量の物差しにならないこと、すなわち配分戦略の比較で見た busy 膨張の主因が非対称コアであることが、推測ではなく測定値で言えるようになった。

\subsection{主張 2 は二つに分けて判定する}

swap 説は棄却できる。
グラフ 1 コピーのメモリ使用量を実測すると約 151\,MB で、worker 8 個と親を合わせても約 1.4\,GB と 16\,GB に十分収まる。

一方で共有資源の競合自体は実在する。
P コアだけの w=4 でも効率が 79\% に落ち、E コアだけの w=4 でも 73\% に落ちる。
idle はどちらもほぼ 0\% なので、これは待ちではなく「同じ種類のコアを複数同時に使うと 1 コアあたりの速度が下がる」現象である。
候補はメモリ帯域・キャッシュの共有と、複数コア稼働時のクロック引き下げで、この二つの内訳は未検証である。

\subsection{環境の影響も定量できた}

Chrome を閉じただけで interleave w=8 の compute は 24.8 秒から 18.2 秒へ約 26\% 縮んだ。
配分戦略の比較は全条件を同じ環境で測っているので相対比較の結論は揺るがないが、絶対値の比較は同一セッションでも背景アプリの状態まで揃える必要がある。
測定前に \texttt{ps} で環境を記録する手順をプロトコルに加えたのはこの観測による。

クリーンな環境での interleave w=8 は compute 18.2 秒（対 w=1 で 5.3 倍、total でも 4.9 倍）となり、このマシンでの到達点が更新された。
P コア換算 5.26 個分という値は、E コアの下限速度比 0.26 から作る単純モデル（$4 + 4\times0.26 \approx 5.1$）と整合する。
ただしこの一致は次節で見直す。

\section{E コアの速度の正体}
\label{sec:ecore-speed}

\subsection{意図}

切り分け実験で得た速度比 0.26 は、\texttt{taskpolicy -b} が優先度も同時に下げるため、「約 4 倍遅い」が E コアの実力なのかクロック制限の結果なのか区別できない。
この疑問（クロック比が機械的に出ているだけではないか）を検証する。

\subsection{純計算ループによる切り分け}

メモリをほぼ使わない整数演算ループを同じ方法で比較すると（3 回の最小値）、P コア 0.64 秒に対し E コア固定で 3.68 秒、速度比 5.75 倍だった。
soc の実測（3.79 倍）より差が大きい。
クリーク列挙にはメモリ待ちが含まれ、メモリ待ちの時間はクロックを落としても伸びないため、低クロックの被害が薄まる。
つまりこの差は「クロック制限が主因、メモリ待ちが緩和要因」という構図を支持する。
副産物として、この対比はクリーク列挙が計算律速寄りであることの傍証にもなっている。

\subsection{公表情報との突き合わせ}

Apple はスペックページにクロック周波数を載せない方針のため、公式値は存在しない。
代わりに、powermetrics による第三者測定（The Eclectic Light Company、Oakley 2024）がある。
M4 の E コアは 1{,}020〜2{,}592\,MHz の 7 段階 DVFS で動作し、background QoS のタスクは約 1{,}050\,MHz（最低周波数付近）に張り付くと報告されている。

これで実測値が分解できる。
P コア約 4.4\,GHz 対 1.05\,GHz のクロック比は 4.2 倍で、純計算ループの実測 5.75 倍との残差（$\times1.37$）が、E コアの IPC（1 クロックあたりの仕事量）が P コアより低い分と読める。
すると全力（2.59\,GHz）の E コアの速度比は $0.59 \div 1.37 \approx 0.43$ と推定できる。
これは直接測定ではなく、第三者測定のクロック値と自前の IPC 係数からの推定である。

\subsection{推定値でモデルを組み直す}

この $r \approx 0.43$ で w=8 をモデル化し直すと、理想のスループットは $4 + 4\times0.43 \approx 5.7$ P コア分になる。
実測は 5.26 なので効率 92\%、残り約 8\% が共有資源の競合ロスという内訳になる。
前節で 0.26 を使ったモデル（$4 + 4\times0.26 \approx 5.1$）が実測と合ったのは、E コアを遅く見積もる誤差と競合ロスを無視する誤差が打ち消し合った偽の一致だった。
いずれのモデルでも、配り方の改善ではこれ以上縮まないという結論は変わらない。

\subsection{方法論上の発見}

測定の過程で、E コア速度の測り方と熱ドリフトについて二つの知見を得た。

\begin{itemize}
  \item macOS ではユーザー空間の工夫でプロセスを「全力の E コア」に固定する手段が事実上ない。\texttt{taskpolicy -c utility} は P コアが空いていると E に固定されず（P と同速）、P コアを飽和させて押し出す方法もスケジューラがプロセスをコア間で回すため固定できない（1.19 倍しか落ちない）。確定させるには \texttt{sudo} で powermetrics を使い実クロックを観測するしかない（未実施、将来課題）。
  \item 長時間のベンチマーク連続実行でマシン自体が遅くなる。約 1.5 時間の測定後、同じソロの純計算ループが 0.64 秒から 0.85 秒へ 32\% 劣化した（熱によるクロック低下と見られるが pmset に記録はなく未確認）。同一セッション内でも時間が離れた測定の絶対値比較にはこのドリフトが乗る。比較したい条件は時間的に隣接させて測ることをプロトコルに加えた。
\end{itemize}

\section{共有メモリアプローチ}
\label{sec:shm}

\subsection{意図}

フェーズ別分解で、並列化の固定費が startup（グラフの \texttt{pickle} 転送）に集中することがわかった。
この固定費自体を消せば、フォーク済みの 6 グラフのように並列化が損だったケースを得に転じられる。
起動方式を変える二つのアプローチを試す。

\begin{itemize}
  \item \textbf{fork}：\texttt{multiprocessing.get\_context("fork")} で worker を起動する。子プロセスは親のメモリを copy-on-write で継承するため、グラフの \texttt{pickle} 転送も worker 側の再構築も発生しない。バッチの作り方は block と同一にし、block との差がすべて起動方式に帰着するようにした。macOS の既定が spawn なのは、スレッドを持つプロセスを fork すると子がデッドロックしうるためだが、本実装は fork 時点で単一スレッドの純 Python なので安全である。
  \item \textbf{shm}：起動方式は spawn のまま、グラフを CSR 形式の int64 配列 4 本（頂点 ID・indptr・隣接リスト・縮退順序）として \texttt{multiprocessing.shared\_memory} に置き、worker が名前で attach する。再帰カーネルは Python の set 演算に依存し、set は共有メモリに置けないため、worker は共有配列から自分用の dict-of-sets を再構築する。消えるのは \texttt{pickle} 化とパイプ転送だけで、worker 側の構造構築は残る。
\end{itemize}

spawn の startup を「interpreter 起動＋転送＋構築」に分解したとき、転送だけを消すのが shm、三つとも消すのが fork という対比になる。

\subsection{方法}

測定条件はバッチ 64、workers=1,4,8、各条件 3 回の最小値である。
同一セッションで fork、shm の順に連続実行した。
測定前に \texttt{ps} を確認し、前セッションの残骸である busy loop 2 プロセス（background QoS で E コアを占有）を停止してから開始した。
極大クリーク数は soc の 22{,}226{,}173 ほか全条件で一致する。
倍率は同一手法内の対 w=1 の total 比である。

\subsection{結果}

9 グラフの全数フェーズ別内訳は付録に回し、ここでは各グラフの w=4,8 での倍率だけを表~\ref{tab:shm-summary} にまとめる。
倍率が 1 を超えれば並列化の得である。

\todo{付録に全数表}

\begin{table}[htbp]
  \centering
  \caption{fork と shm の 9 グラフ要約（対 w=1 total 倍率、3 回の最小値）}
  \label{tab:shm-summary}
  \begin{tabular}{lrrrr}
    \toprule
    & \multicolumn{2}{c}{fork} & \multicolumn{2}{c}{shm} \\
    \cmidrule(lr){2-3}\cmidrule(lr){4-5}
    グラフ & w=4 & w=8 & w=4 & w=8 \\
    \midrule
    木 $n{=}20$k          & 1.97 & 1.88 & 1.22 & 0.97 \\
    2 次元格子 $100\times100$ & 2.01 & 1.92 & 1.18 & 0.86 \\
    疎ER $n{=}20$k        & 2.09 & 2.28 & 1.47 & 1.17 \\
    BA hubs $n{=}20$k     & 2.35 & 2.61 & 1.58 & 1.26 \\
    密ER $n{=}120$        & 1.08 & 1.07 & 1.07 & 0.99 \\
    Moon–Moser $k{=}11$   & 0.97 & 0.95 & 0.96 & 0.85 \\
    ca-AstroPh            & 3.09 & 3.91 & 2.35 & 2.20 \\
    ca-HepPh              & 2.76 & 2.86 & 2.23 & 1.80 \\
    soc-sign-epinions     & 2.65 & 3.05 & 2.78 & 3.24 \\
    \bottomrule
  \end{tabular}
\end{table}

\subsection{fork は startup をほぼ消した}

全 9 グラフで fork の startup は 0.002〜0.031 秒に収まった。
spawn では worker 数に比例して伸びていた固定費（soc で w=8 の 1.273 秒）が、fork では w=8 でも 0.031 秒である。
この結果、spawn では並列化が全敗だった小さい 6 グラフのうち 4 つ（木・格子・疎ER・BA）が w=4 で約 2 倍の得に転じた。
並列化の損益分岐点は「compute 数秒」から「compute 数十ミリ秒」まで下がったことになる。
残る密ER と Moon–Moser が改善しないのは startup のせいではなく、compute 自体が縮まないためである。
単一の重い部分木が支配的で、バッチ並列の粒度では割れない。

\subsection{shm は転送しか消せない}

shm の startup は spawn より小さいが fork より一桁大きい（soc w=1 で 0.413 秒、小グラフで 0.04〜0.16 秒）。
残るのは親側の CSR パック、spawn の interpreter 起動、worker 側の set 再構築である。
この実験は spawn の startup の内訳を分離したことになる。
転送を消すだけでは足りず、構造構築と interpreter 起動が残る限り、小グラフの損益分岐点は大きくは動かない。
小さい 6 グラフでの得は最大 1.58 倍（BA w=4）にとどまり、fork に全面的に劣る。

\subsection{shm の compute が fork より速い}

予想外の観測として、shm の compute が fork より一貫して速かった。
soc w=1 で 85.2 秒対 96.2 秒（11\%）、ca-AstroPh w=1 で 0.703 秒対 0.841 秒である。
再帰カーネルは同一なので、手法のコード差では説明できない。

有力な機構は fork の copy-on-write の遅延コストである。
Python はオブジェクトを読むだけでも参照カウントを書き換えるため、fork で共有したページは compute 中に次々とコピーされる。
つまり fork の startup の安さの一部は、コピーを compute に付け替えたものと読める。
一方 shm の worker は自分で作ったコンパクトな set を読むのでこのコストがない。
熱ドリフトなら後に走った shm が遅くなるはずで方向が逆なので、環境説より CoW 遅延コスト説が有力である。
ただし切り分け（例：\texttt{gc.freeze()} を使った fork の再測定）は未検証である。

\subsection{到達点}

soc の w=8 total は spawn 37.8 秒（前回測定）、fork 31.7 秒、shm 26.5 秒と縮んだ。
ただし spawn の値は別セッションの測定で、絶対値は環境に極めて敏感（Chrome 開閉で 26\%）なので、厳密な手法間比較として読んでよいのは同一セッションで隣接して測った fork 対 shm のみである。

結論として、startup 固定費を消す目的には fork が最小の変更で最大の効果を与える。
shm は startup 削減では fork に劣るが、compute を含めた total では soc で最速だった。
両者の両立（fork と共有配列の併用、あるいは \texttt{gc.freeze}）は将来課題である。

\chapter{おわりに}
\label{ch:conclusion}

\section{本研究のまとめ}

本研究は、極大クリーク列挙問題における Eppstein アルゴリズムの並列化を、単一の実装で終わらせず、測定の連鎖として進めてきた。
ここでその連鎖を振り返る。

出発点は逐次アルゴリズム同士の比較である。
十種類のグラフのうち、実データで最も重い soc-sign-epinions では、Tomita らのピボット版 CLIQUES が167.3秒であるのに対し、縮退順序を用いる Eppstein は94.9秒だった。
疎なグラフほど Eppstein が有利という傾向は、縮退数 $d$ とグラフサイズ $n$ の比 $d/n$ でおおむね説明できた。
この結果を受けて、以降の並列化は Eppstein を対象とし、最も重い soc-sign-epinions を主力のベンチマークに定めた。

次に試したのが、最外層の頂点を worker に配るだけの素朴な並列化である。
バッチサイズ512では、worker 数を4以上に増やしても67〜70秒から縮まなかった。
バッチサイズを64に変更しただけで、同じ8 worker が27.7秒まで縮む。
この一つの変更で結果が大きく動いた事実は、並列化の遅さの原因が worker 数ではなくタスクの配り方にあることを示していた。

そこで負荷分布を頂点単位で測定すると、soc-sign-epinions では上位1\%の頂点が全体の94.7\%の時間を占めていた。
偏りの発生源は縮退順序の末尾（疎な外周から剥がしていくため、密なコアが必然的に末尾に残る）である。
一方 ca-HepPh は最も重い1頂点だけで13\%を占め、どう配ってもその1タスクの分割なしにはスピードアップの上限が7.7倍程度に制限される。
どちらも「並列化が効かない」と要約できてしまいそうだが、前者はタスクの配り方で救える偏りであり、後者は部分問題内部を分割しない限り解決しない別種の問題である。

実行時間をロード、縮退順序計算、プロセスプール起動、並列本体の4区間に分解すると、プール起動の固定費が worker 数にほぼ比例して伸びることが分かった（soc で w=1 の0.268秒から w=8 の1.273秒）。
あわせて worker ごとの稼働率を見ると、w=8 では idle が30.4\%に達し、最も忙しい worker と最も暇な worker の差は36.2秒対11.8秒だった。
この idle は、負荷分布の測定で見つけた偏りが worker の遊休として現れたものである。

配分戦略そのものを block、reversed、interleave の三通りで比較すると、interleave は idle を32.6\%からほぼ0.1\%まで解消した。
ところが compute の短縮は4.3\%にとどまり、idle がほぼ消えたにもかかわらず理論上期待される約1.5倍の短縮には遠く及ばなかった。
idle の解消と引き換えに、全 worker の稼働時間の合計は139秒から198秒へ膨張しており、遊休を埋める効果を稼働時間そのものの膨張が打ち消していた。

この膨張の機構を、taskpolicy による P/E コア固定実験で切り分けた。
同じ計算を P コアで実行すると96.0秒、E コアに固定すると363.6秒であり、この条件では E コアは P コアの4分の1程度の速度しか出ない（固定に用いた background QoS はクロックも制限するため、この比は下限値である）。
メモリをほとんど使わない純計算ループでは速度比がさらに開く（5.75倍）ことから、この差の主因はクロック制限であり、クリーク列挙で差が縮むのはメモリ待ちの時間が低クロックの影響を受けないためだと判断できる。
Oakley による powermetrics 測定 \cite{oakley2024} を参照してクロック比から E コアの実力を推定すると、w=8 の実測スループット（P コア5.26個分相当）は、この推定モデル（$4 + 4\times0.43 \approx 5.7$ P コア分、効率92\%）とおおむね整合した。
配分戦略比較で観測した busy 時間の膨張は、この非対称コアが主因だと測定によって裏づけられたことになる。

最後に、並列化のもう一つの固定費であるメモリコピーの削減に取り組んだ。
multiprocessing の起動方式を spawn から fork に変えると、copy-on-write によってグラフの pickle 転送が不要になり、startup は全グラフで0.002〜0.031秒まで縮んだ。
これにより、spawn では並列化がすべて損だった小さいグラフのうち4つが、fork では w=4 で約2倍の得に転じた。
グラフを CSR 形式の配列として shared\_memory に置く shm 方式は、startup の削減幅では fork に劣るものの、compute 自体が fork より一貫して速く、soc-sign-epinions の w=8 では total が26.5秒に達した。
逐次版の94.9秒（別セッションの測定のため参考値）から数えると約3.6倍にあたり、この時点が本研究の到達点である。

\section{得られた知見}

一連の測定から言えるのは、並列化の性能はアルゴリズムの選択だけで決まるのではないということである。
本研究の範囲では、少なくとも四つの要因が性能に効いていた。
プロセスプール起動とデータ転送にかかる\textbf{起動固定費}、縮退順序の末尾に偏る\textbf{負荷分布}、P コアと E コアで速度が大きく異なる（下限実測で約4倍、全力推定でも2倍以上）\textbf{ハードウェアの非対称性}、そして pickle 転送か copy-on-write か共有メモリかという\textbf{メモリ共有方式}である。

これらは独立に効くのではなく、互いを覆い隠す形で作用していた。
配分戦略の改善で負荷分布起因の idle をほぼ消しても、その分の稼働時間が非対称コアの存在によって膨張し、期待した高速化は得られなかった。
逆に言えば、負荷分布だけを見て「配分は改善済みだからこれ以上は望めない」と判断していたら、非対称コアという別の要因を見落としたままになっていた。
本研究で複数の切り分け実験を積み重ねたのは、この種の要因の重なりを一つずつ分離するためである。

ただし、この知見は Apple M4 搭載の1台のマシン、Python の multiprocessing、十種類のグラフという条件のもとで得られたものであり、他のハードウェアや実装、より大規模なグラフでも同じ配分で効くとは限らない。
特に非対称コアの影響は、P コアと E コアを持たない一般的なサーバ環境では前提から成り立たない。
ここで得られた一般化は、あくまで本研究の実測が支える範囲にとどまる。

\section{今後の課題}

本研究の測定の過程で、実施に至らなかった検証や、切り分けが済んでいない機構がいくつか残っている。

\begin{itemize}
  \item \textbf{copy-on-write 説の切り分け}：fork の compute が shm より遅い観測について、CoW の遅延コスト（参照カウント書き換えによるページコピー）が原因だとする説を、\texttt{gc.freeze()} を使った再測定で検証すること。
  \item \textbf{E コアの実クロックの確定}：E コアの速度低下がクロック制限によるという推定を、\texttt{sudo} 権限下の \texttt{powermetrics} による実クロックの直接観測で確定させること。
  \item \textbf{共有資源競合の内訳の切り分け}：P コアのみの並列でも効率が79\%に落ちる競合について、キャッシュ、メモリ帯域、クロック引き下げのどれによるものかを切り分けること。
  \item \textbf{熱ドリフトの機構の確認}：長時間の連続測定で同じ純計算ループが32\%遅くなる現象の原因（クロック低下と見られるが未確認）を、実クロックと温度の記録で確認すること。
  \item \textbf{work stealing とオーバーデコンポジション}：探索木を worker 数よりずっと多いサブツリーに分割し、暇な worker が動的にキューから取る方式を導入し、静的な配分戦略の限界を超えられるかを検証すること。
  \item \textbf{グラフの CSR 化}：shm 方式で導入した CSR 形式による表現を、他の起動方式や、より大きなグラフに対しても一貫して適用すること。
  \item \textbf{ca-HepPh 型の単一重部分問題の内部分割}：縮退順序の1頂点だけが全体時間の大半を占めるグラフに対して、その部分問題自体をさらに分割する方法を検討すること。
  \item \textbf{fork と共有配列の両立}：起動固定費を fork 並みに抑えつつ、shm 並みの compute 速度も得られる構成を検討すること。
\end{itemize}


% 文献は全て英語のため plain を使う。研究室指定があれば jplain 等に差し替える
\bibliographystyle{plain}
\bibliography{refs}

\end{document}
